NIVEL 1
Paso 1. Diseña la arquitectura
Landing │ 
▼ BRONZE │ 
▼ SILVER │ 
▼ GOLD

Bronze: Guardar datos casi idénticos al origen.
Silver: Datos limpios.
Gold: Métricas para negocio.

Paso 2. - Bronze
R1.1
Debes cargar todas las fuentes:
Colombia: JSON CDC
México: CSV
CRM : CSV
Catálogos: JSON y CSV

--Columnas:
fecha_ingesta
archivo_origen
pais_fuente
timestamp_procesamiento

--Justificación en DECISIONES.md
Se eligió parquet para las tablas bronze por ser columnar, comprimido, ampliamente soportado por Spark y compatible con una futura migración a Delta Lake en Databricks.

Paso 3 - Hallazgos de calidad
R1.2 -NO debes solo limpiar.

Debes cuantificar.
Hallazgo	Fuente	Cantidad	%	Impacto	Decisión
Canal desconocido	MX	15	0.03%	No clasifica negocio	Excluir
Comisión nula	CO	20	0.02%	Afecta ingresos	Cuarentena
Comercio inexistente	CRM	35	0.04%	No puede enriquecer	Mantener con UNKNOWN



------------------------
## Olas de entrega
Procesa cada ola como una corrida independiente del pipeline, en orden: `ola_1` -> `ola_2`.

| Ola | Contenido (por fecha de llegada a la landing) |
|-----|------------------------------------------------|
| ola_1 | Archivos que llegaron entre 2026-06-01 y 2026-06-30 |
| ola_2 | Archivos que llegaron entre 2026-07-01 y 2026-08-02 |


## Definiciones de negocio (vigentes)
- **Fecha de negocio**: fecha local de la transaccion en su pais (CO: America/Bogota, MX: America/Mexico_City).
- **Transaccion efectiva**: estado final APROBADA. Un reverso posterior (REVERSADA/CANCELADA) anula la transaccion en su fecha de negocio original.
- **Ingreso**: comision de la transaccion efectiva en moneda local.
- Transacciones de comercios o productos de prueba **no** hacen parte de las cifras oficiales.
- Montos por fuera de los limites del catalogo del producto no son validos para cifras oficiales.